Ich möchte Ihnen zeigen, wie sich die Darstellungstheorie -- insbesondere also die Darstellung der Gruppen -- auffassen lässt als eine Theorie der Modul- und Idealklassen; und wie diese letzteren sich wieder einem verallgemeinerten Gruppenbegriff unterordnen, den Gruppen mit Operatoren.\footnote{Die Gruppen mit Operatoren gehen auf W. Krull (Math. Zeitschr. 23) und O. Schmidt (Math. Zeitschr. 29) zurück. Auch die hier gebrachte Behandlung des Problems der Reduktion einer gegebenen Darstellung rührt -- in der Beschränkung auf kommutative Darstellungskörper -- von W. Krull her (Theorie und Anwendung der verallgemeinerten Abelschen Gruppen, Heidelberger Ber. 1926).}
